%% base_ieee.tex

\documentclass[journal]{IEEEtran}
%

\usepackage{graphicx}
\usepackage{float}
\usepackage{subfig}
\usepackage{cite}


% *** Do not adjust lengths that control margins, column widths, etc. ***
% *** Do not use packages that alter fonts (such as pslatex).         ***
% There should be no need to do such things with IEEEtran.cls V1.6 and later.
% (Unless specifically asked to do so by the journal or conference you plan
% to submit to, of course. )


% correct bad hyphenation here
\hyphenation{op-tical net-works semi-conduc-tor}


\begin{document}
%
% paper title
% Titles are generally capitalized except for words such as a, an, and, as,
% at, but, by, for, in, nor, of, on, or, the, to and up, which are usually
% not capitalized unless they are the first or last word of the title.
% Linebreaks \\ can be used within to get better formatting as desired.
% Do not put math or special symbols in the title.
\title{Frequency Control for Inverter-based Resources}
%
%
% author names and IEEE memberships
% note positions of commas and nonbreaking spaces ( ~ ) LaTeX will not break
% a structure at a ~ so this keeps an author's name from being broken across
% two lines.
% use \thanks{} to gain access to the first footnote area
% a separate \thanks must be used for each paragraph as LaTeX2e's \thanks
% was not built to handle multiple paragraphs
%

\author{Sean~Jennings}%
        % ,~\IEEEmembership{Member,~IEEE,}



% The paper headers
% \markboth{Journal of \LaTeX\ Class Files,~Vol.~14, No.~8, August~2015}%
% {Shell \MakeLowercase{\textit{et al.}}: Bare Demo of IEEEtran.cls for IEEE Journals}
% The only time the second header will appear is for the odd numbered pages
% after the title page when using the twoside option.
% 
% *** Note that you probably will NOT want to include the author's ***
% *** name in the headers of peer review papers.                   ***
% You can use \ifCLASSOPTIONpeerreview for conditional compilation here if
% you desire.



\maketitle

% As a general rule, do not put math, special symbols or citations
% in the abstract or keywords.
\begin{abstract}
The abstract goes here.

This is a mf abstract!
\end{abstract}

% Note that keywords are not normally used for peerreview papers.
% \begin{IEEEkeywords}
% IEEE, IEEEtran, journal, \LaTeX, paper, template.
% \end{IEEEkeywords}






% For peer review papers, you can put extra information on the cover
% page as needed:
% \ifCLASSOPTIONpeerreview
% \begin{center} \bfseries EDICS Category: 3-BBND \end{center}
% \fi
%
% For peerreview papers, this IEEEtran command inserts a page break and
% creates the second title. It will be ignored for other modes.
\IEEEpeerreviewmaketitle



\section{Introduction}
\IEEEPARstart{T}{he} penetration of renewable energy sources (RES)
projected to increase over the coming decades, with gains highly powered by
solar and wind resources. The NREL report %\ref
has solar and wind energy increasing to >50\% share by 2050. 
Solar photovoltaics (PV) and Type III and Type IV wind turbine generators
(WTG) are both inverter-based resources (IBRs), as opposed to synchronous
generators (SGs) of which the conventional power grid is mostly composed.

Throughout the history of the power grid, reliability has been a prime concern.
A core aspect of power system reliability is frequency control, the ability
to maintain the AC voltage and current signals close to a given setpoint
(e.g. 60 Hz in the USA). If frequency deviates too far from its setpoint,
protection mechanisms of the grid, such as under-frequency load shedding
(UFLS) are triggered to prevent damage to equipment and to arrest further
instability.  %\ref?
Load shedding means that some customer(s) have lost power, a serious grid
reliability issue.

When the grid is near equilibrium, its frequency is closely related
to the real power transfer between its nodes. When a system fault, such as a
generator going offline, causes a power imbalance in the system, 
the rest of the generators increase power output to balance
generation and load. To transfer power between nodes requires the phase angle
between nodes in the system to shift, which results in transient frequency
deviations.

The rate at which frequency changes (RoCoF) is related to underlying physical
system from which the energy is derived. Conventionally,
thermal power plants use SGs to convert the underlying mechanical energy
contained in steam into electrical energy. The RoCoF is limited by
the mechanical {\bf inertia} of the SG as governed by the swing equation.

On the other hand, IBRs connect an energy resource such as a PV cell,
WTG, or a battery energy storage system (bESS) to the grid via a 
power electronic (PE) inverter which converts a DC source to an AC output.
The dynamics governing the energy resource as well as the control algorithm.
Because the IBR is not directly connected to any form of mechanical inertia,
IBR do not have inertia and that systems with high penetrations of IBRs
are {\it low-inertia} systems. \cite[text]{}

The control algorithms used by IBRs can be broadly divided into two strategies:
\begin{enumerate}
  \item Grid-following Inverters (GFL), and
  \item Grid-forming Inverters (GFM).
\end{enumerate}

This distinction, however, is informal. To date, there is no IEEE standard
on the definition of GFM. NREL defines a GFM to be any inverter that does 
not operate with a phase-lock loop (PLL). %ref??

There has been some debate on the security of IBRs. %ref{foundation-and-challenges}
posits that "low inertia likely implies low security" since
lack of mechanical inertia means that the stability of the system
relies on the software implementation of control algorithms. Thus,
there is a growing body of research aimed at demonstrating the 
robustness of the control strategies deployed on IBRs.

While the deployment of increasing penetrations of IBR to the large-scale
grid may not be gauranteed, it remains one of the currently most 
economically feasible paths to net-zero carbon emissions. Furthermore,
any potential deployment will necessarily be gradual process: the entirety of
SG-based generation cannot be replaced with generation based on IBRs
overnight. Therefore, it is of vital importance to study the 
dynamics of IBR-based frequency control in the variety of contexts
anticipated in the forthcoming energy transition. Since this
entails nothing less than a complete rearchitecting of the energy
grid, there are numerous open research questions in this field
which will be enumerated.

This literature review focuses on the challenges and oppurtunities along
the path to incorporating high levels of IBR penetration into the grid 
at a large scale. It is organized as follows: section \ref{sec:math-models}
summarizes the mathematical formulation for frequency control; section
\ref{sec:control-strategies} enumerates and classifies the many
of control schemes that have been proposed in literature; section
\ref{sec:commericial-demonstrations} summarizes the major commericial
projects which have been deployed using GFM technology to-date.

\section{A Mathematical Formulation for Frequency Control}
\label{sec:math-models}

\subsection{Synchronous Generators}
 We begin our summary of the literature by reviewing the dynamics of SGs, 
 which is a topic that has been well-established
 for decades. Despite it's lack of novelty, the topic remains highly
 relevant to the problem of IBR frequency control since
 \begin{enumerate}
  \item a substantial portion of the prevailing industry terminology
    is based around these terms, and
  \item SG dynamics interact with IBR dynamics via their network
    interconnections within the grid.
 \end{enumerate}

%  A SG is composed of three major components from a control standpoint,
% the excitation system, the 

 The relationship between power and frequency in a SG is governed by
 \begin{equation}
  J\dot{\omega} = p_m - p_e
  \label{eq:swing-equation}
 \end{equation}
 where $J$ is the system intertia, $p_m$ is the mechanical power input to
 the generator, and $p_e$ is the electrical power which is absorbed
 by connected loads.

Since a change in frequency signals there is a power imbalance within the
system, the rotor speed of an SG is typically regulated by a
feedback-mechanism called the governor. This feedback couples a change in 
frequency, $\omega$, to a (typically proportional) change in the power
delivered by the prime mover $p_m$:
\begin{equation}
  \dot{p_m} = - \frac{\omega_{set}}{R} * (\omega - \omega_{set}) + p_{m,set}
  \label{eq:governor-control}
\end{equation}

Combining equation \ref{eq:swing-equation} and \ref{eq:governor-control}
gives a second order relationship between frequency, $\omega$ and 
mechanical power delivered, $p_m$.
 
\subsection{IBR Grid Integration}
In contrast, the frequency-real power relation for IBR resources are not
coupled by mechanical inertia of the swing equation,
but rather are determined by the underlying energy source for the IBR system.
%\ref{interactivity}
defines makes the assumption that power and energy is available from a bESS,
but the same conclusions hold if the energy source is curtailed PV or wind power.
This assumption is a common and vital one: it informs the design of
the DC side of an IBR. %\ref{battery-sizing}
elaborates on the relationship between the control parameters and
the power requirements for the bESS.

%\ref{interactivity}
demonstrates that the timescale of IBR reactions are an order of magnitude
shorter than the mechanical timescale (~1 second) which SGs operate.
Thus, when considering the interaction between the two we find that 
% $p_{m,G} \approx p_{m,Go}$ is a
on this short timescale, the IBR response is first-order with respect to the
electrical power $p_{m,I}$ drawn from the circuit:
\begin{equation}
  \dot{p_{m,I}} = \frac{2 \pi (p_{e,I} - p_{m,I})}{\tau_I}
\end{equation}

It is well-documented that if the grid contains above a certain penetration
of GFLs, instability occurs. Therefore, for grids with high
IBR penetrations, a certain percentage of GFMs are necessary to maintain
stability. However, GFMs can also have their own drawbacks: they are
require more specialized hardware and control algorithms so are
therefore more expensive.

It has been shown that the placement of GFMs within a power system
topology has an important impact on overall system stability in a 
system with mixed GFL, SG, and GFM-based resources. However,
the exact interaction between GFLs and GFMs remains an open research
question. In particular, the ratio of GFLs to GFMs at which instability
arises and the optimal placement of GFMs within a given power system
remain under-studied questions.

A major difference between the physics of IBRs versus SGs is that
typically IBRs have a much lower current threshold that they
can produce based on the limits of the underlying PEs. As a result,
they have a smaller short-circuit ratio, which corresponds to 
a lower system strength. In general, a weaker system strength
implies a decrease in the voltage stiffness, which can impair
IBR performance. GFMs can provide "voltage strength" which is another
advantage over GFLs. Since this review is focused on frequency
control of IBRs, we will consider further discussion on the topic
as out of scope and will conclude with the note that voltage and
frequency control are inherently coupled and modeling these 
interactions is vital for demonstrating system stability.

\section{GFM Control Strategies}
\label{sec:control-strategies}

Literature on GFM control has existed since the mid-1990s %\ref{?}
since it is a technical requirement for microgrids (MGs), which may
not have SGs and therefore must establish their own frequency.
(move to Intro?)

A number of schemes for controlling GFMs have been proposed, including
\begin{enumerate}
  \item droop control
  \item multi-loop droop control
  \item virtual synchronous machine (VSM)
  \item virtual oscillating control (VOC)
\end{enumerate}

\subsection{Virtual Inertia-based Schemes}
Both VSM and VOC strategies aim to emulate SGs by designing controllers
which exhibit similar dynamical properties to an SG. The degree
to which it is possible to fully emulate an SG is again driven by 
the physics of the underlying system, but the studies have shown
the capability of PE to emulate a SG within both simulation %\ref{virtual-inertia-cs}
and on physical devices.

These methods ensure that the RoCoF after a distrubance remains constrained
within a certain limit. This is important because various components of
power system protection make assumptions on RoCoF, so that changes
in the dynamical properties of underlying generation may invalidate
protection-related assumptions.


\subsection{Droop Control}
To the best of the author's knowledge, there exists no formal definition
of "droop control" within GFM literature. The origins of the
term "droop" is vestigial terminology from the tradition of SGs, where
it presumably refers to the "droop" of the mechanical centrifugal governor
control mechanism [add figure].

The commonality between droop control techniques is the inverse
relationship between frequency and power generation, such that
frequency is essentially used as a signal within the grid that a power
imbalance has occurred and that other generation sources should
increase their power output accordingly.


\subsection{Optimization-based techniques}
[TODO: foundations and challenges paper has interesting notes on optimal
control parameter tuning, etc., read this more thoroughly!]



\section{Commercial Demonstrations of GFM systems}
\label{sec:commericial-demonstrations}

This section analyzes the scientific documentation of the
current state of the GFM industry's commercial development,
with the objective of illuminating the degree to which
the control strategies and system models of sections
\ref{sec:math-models} and \ref{sec:control-strategies}
have been tested and validated on physical systems rather
than simply simulation.

\subsection{Distributed IBRs}
The IBRs which have been deployed to-date within PV and WTG
have been primarily GFL technology. (Is this true for WTG??)
Since converting a GFL IBR to a GFM IBR requires upgrading
the hardware itself, along with a high capital cost. The projected
increase in GFL penetration makes the understanding the relationship
between GFL and GFM dynamics an even more pressing topic.

\subsection{Microgrids}
The Maui power system is a commonly used example of a MG which has 
high instantaneous power of IBRs. Other small-scale  grids
(e.g. grids on small islands) have demonstrated the technical
feasibility of GFM technology while simultaneously presenting 
discrepancies from existing simulation results.

\subsection{Utility/Grid-scale Implementations}
The authors of [ref]%\ref{}
provide a overview of grid-scale systems which have implemented GFM
technology. The initial results have been promising and each provides
unique insight into what the future of high penetration integration of
IBRs might look like.



\section{Conclusion}
Given the projected rapid growth in IBRs, the instability of a
zero-inertia system based entirely off of GFL technology, as well
as the current rapid deployment of GFL technology, there is a clear
necessity for standards and commericalization of GFM technoogy.

A prerequisite for such standards is thorough characterization
of the stability characteristics of hybrid power systems. The
large body of active research into GFM control algorithms
has provided many promising results for the technical feasibility.

Still, many questions remain. The author is particularly interested
in further investigating the relationship between PS topology, GFM/GFL
hybrid dynamics, as well as a more detailed understanding of the
convergence properties of various optimization-based control algorithms.


Hello, bye


% An example of a floating figure using the graphicx package.
% Note that \label must occur AFTER (or within) \caption.
% For figures, \caption should occur after the \includegraphics.
% Note that IEEEtran v1.7 and later has special internal code that
% is designed to preserve the operation of \label within \caption
% even when the captionsoff option is in effect. However, because
% of issues like this, it may be the safest practice to put all your
% \label just after \caption rather than within \caption{}.
%
% Reminder: the "draftcls" or "draftclsnofoot", not "draft", class
% option should be used if it is desired that the figures are to be
% displayed while in draft mode.
%
%\begin{figure}[!t]
%\centering
%\includegraphics[width=2.5in]{myfigure}
% where an .eps filename suffix will be assumed under latex, 
% and a .pdf suffix will be assumed for pdflatex; or what has been declared
% via \DeclareGraphicsExtensions.
%\caption{Simulation results for the network.}
%\label{fig_sim}
%\end{figure}

% Note that the IEEE typically puts floats only at the top, even when this
% results in a large percentage of a column being occupied by floats.


% An example of a double column floating figure using two subfigures.
% (The subfig.sty package must be loaded for this to work.)
% The subfigure \label commands are set within each subfloat command,
% and the \label for the overall figure must come after \caption.
% \hfil is used as a separator to get equal spacing.
% Watch out that the combined width of all the subfigures on a 
% line do not exceed the text width or a line break will occur.

% \begin{figure*}[!t]
% \centering
% \subfloat[Case I]{\includegraphics[width=2.5in]{box}%
% \label{fig_first_case}}
% \hfil
% \subfloat[Case II]{\includegraphics[width=2.5in]{box}%
% \label{fig_second_case}}
% \caption{Simulation results for the network.}
% \label{fig_sim}
% \end{figure*}
%
% Note that often IEEE papers with subfigures do not employ subfigure
% captions (using the optional argument to \subfloat[]), but instead will
% reference/describe all of them (a), (b), etc., within the main caption.
% Be aware that for subfig.sty to generate the (a), (b), etc., subfigure
% labels, the optional argument to \subfloat must be present. If a
% subcaption is not desired, just leave its contents blank,
% e.g., \subfloat[].


% An example of a floating table. Note that, for IEEE style tables, the
% \caption command should come BEFORE the table and, given that table
% captions serve much like titles, are usually capitalized except for words
% such as a, an, and, as, at, but, by, for, in, nor, of, on, or, the, to
% and up, which are usually not capitalized unless they are the first or
% last word of the caption. Table text will default to \footnotesize as
% the IEEE normally uses this smaller font for tables.
% The \label must come after \caption as always.
%
%\begin{table}[!t]
%% increase table row spacing, adjust to taste
%\renewcommand{\arraystretch}{1.3}
% if using array.sty, it might be a good idea to tweak the value of
% \extrarowheight as needed to properly center the text within the cells
%\caption{An Example of a Table}
%\label{table_example}
%\centering
%% Some packages, such as MDW tools, offer better commands for making tables
%% than the plain LaTeX2e tabular which is used here.
%\begin{tabular}{|c||c|}
%\hline
%One & Two\\
%\hline
%Three & Four\\
%\hline
%\end{tabular}
%\end{table}


% Note that the IEEE does not put floats in the very first column
% - or typically anywhere on the first page for that matter. Also,
% in-text middle ("here") positioning is typically not used, but it
% is allowed and encouraged for Computer Society conferences (but
% not Computer Society journals). Most IEEE journals/conferences use
% top floats exclusively. 
% Note that, LaTeX2e, unlike IEEE journals/conferences, places
% footnotes above bottom floats. This can be corrected via the
% \fnbelowfloat command of the stfloats package.








% if have a single appendix:
%\appendix[Proof of the Zonklar Equations]
% or
%\appendix  % for no appendix heading
% do not use \section anymore after \appendix, only \section*
% is possibly needed

% use appendices with more than one appendix
% then use \section to start each appendix
% you must declare a \section before using any
% \subsection or using \label (\appendices by itself
% starts a section numbered zero.)
%


% \appendices
% \section{Proof of the First Zonklar Equation}
% Appendix one text goes here.

% % you can choose not to have a title for an appendix
% % if you want by leaving the argument blank
% \section{}
% Appendix two text goes here.
% We appending to the appendix.


% use section* for acknowledgment
% \section*{Acknowledgment}


% The authors would like to thank...


% Can use something like this to put references on a page
% by themselves when using endfloat and the captionsoff option.
\ifCLASSOPTIONcaptionsoff
  \newpage
\fi



% trigger a \newpage just before the given reference
% number - used to balance the columns on the last page
% adjust value as needed - may need to be readjusted if
% the document is modified later
%\IEEEtriggeratref{8}
% The "triggered" command can be changed if desired:
%\IEEEtriggercmd{\enlargethispage{-5in}}

% references section

% can use a bibliography generated by BibTeX as a .bbl file
% BibTeX documentation can be easily obtained at:
% http://mirror.ctan.org/biblio/bibtex/contrib/doc/
% The IEEEtran BibTeX style support page is at:
% http://www.michaelshell.org/tex/ieeetran/bibtex/
\bibliographystyle{IEEEtran}
% argument is your BibTeX string definitions and bibliography database(s)
\bibliography{IEEEabrv}
% \begin{thebibliography}{1}
%   \bibitem{IEEEhowto:kopka}
%   hello
%   % \bibitem
% \end{thebibliography}
%
% <OR> manually copy in the resultant .bbl file
% set second argument of \begin to the number of references
% (used to reserve space for the reference number labels box)
% \begin{thebibliography}{1}

% \bibitem{IEEEhowto:kopka}
% H.~Kopka and P.~W. Daly, \emph{A Guide to \LaTeX}, 3rd~ed.\hskip 1em plus
%   0.5em minus 0.4em\relax Harlow, England: Addison-Wesley, 1999.

% \end{thebibliography}



% that's all folks
\end{document}


